% =====================================================================
%  STATEMENT DOCUMENT.  Informal conjecture upfront, then two sections
%  -- The setting, The conjecture -- following exactly the structure
%  and environment conventions of the sibling k-cr and groth16
%  statements (Section -> prose setting -> formal
%  \definition/\conjecture/\remark in the conjecture section).  No
%  Bibliography section: the original document cited no external
%  source, and none is added here.
%
%  This document originally stated TWO conjectures side by side,
%  "Secret seed" and "Public seed".  The public-seed one is no longer
%  open: resolution.tex in this folder proves it (Theorem + Corollary,
%  constant c=2), and it is published, proved, on the live site as
%  c/0004 (category: research-solved).  The secret-seed conjecture
%  remains open and unattempted, published as c/0005.  Per the
%  author's instruction, this statement keeps only the still-open
%  secret-seed conjecture; the public-seed bound is mentioned in
%  Remark~\ref{rem:pub} purely as resolved context, not restated as a
%  conjecture of its own.
%
%  Built on the shared conjura-conjecture.cls house style -- see
%  latex/conjectures/_template/.
% =====================================================================
\documentclass{conjura-conjecture}

\newcommand{\Fun}{\mathsf{Fun}}
\newcommand{\Adv}{\mathbf{Adv}}
\newcommand{\Gm}{\mathsf{Game}}
\newcommand{\Pred}{\mathrm{Pred}}
\newcommand{\Ext}{\mathrm{Ext}}
\newcommand{\pred}{\mathrm{pred}}
\newcommand{\ext}{\mathrm{ext}}
\newcommand{\extpub}{\textup{ext-pub}}
\newcommand{\sd}{\mathit{sd}}
% overrides the class's plain-arrow \sample with the "sampled from" dollar-arrow this note needs
\renewcommand{\sample}{\stackrel{{\scriptscriptstyle\$}}{\leftarrow}}
\newcommand{\SD}{\mathrm{SD}}
\newcommand{\calK}{\mathcal{K}}
\newcommand{\calD}{\mathcal{D}}
\newcommand{\calR}{\mathcal{R}}
\newcommand{\ret}{\textbf{return}\ }

\runninghead{LHL FOR RANDOM-ORACLE SOURCES}

\begin{document}

\cjkicker{CONJURA \textperiodcentered{} OPEN PROBLEM}
\cjtitle{Randomness Extraction from Unpredictable Random-Oracle Sources}
\cjsubtitle{A Leftover-Hash-Lemma Conjecture for the Secret-Seed Game}
\cjstatus{Statement: human-written, not yet formalized. Proof: open -- no attempt yet.}
\cjcategory{Information-Theoretic Cryptography.}
\cjkeywords{leftover hash lemma; randomness extraction; unpredictable sources; random-oracle
  model; universal hashing; seed-dependent extractors.}

\vspace{0.4em}

\begin{informalconjecture}
For an extractor built from a random function with a $K$-element seed space, and any source
whose output no unbounded predictor can guess with probability better than $\epsilon$ -- even
given the extractor's entire function table -- the extracted value should look uniform to an
unbounded distinguisher that does not see the seed, up to statistical distance about
$\sqrt{(\epsilon R + \log D)/K}$. This is the shape of the classical leftover hash lemma, but
with unpredictability against an omniscient adversary in place of min-entropy as the source
restriction, and a single universal constant working for every choice of parameters.
\end{informalconjecture}

\section{The setting}\label{sec:setting}

An extractor built from a random function $H$ maps a seed and an unpredictable source's output
to a value that should look uniform to an unbounded distinguisher holding the source's auxiliary
output and the oracle's entire table.  Unpredictability is measured against an equally unbounded
predictor with the same information.  Whether the distinguisher also learns the seed splits the
question into two games; this note concerns the harder one, where the seed stays hidden from it.
Definition~\ref{def:games} formalizes both notions.

\smallskip
\noindent\textbf{Notation.}
Fix nonempty finite sets $\calK$ (seeds), $\calD$ (inputs), $\calR$ (outputs), and write
$K := |\calK|$, $D := |\calD|$, $R := |\calR|$.  Let $\Fun(\calK \times \calD, \calR)$ be the set
of all functions $\calK \times \calD \to \calR$, let $\SD(\cdot,\cdot)$ denote statistical
distance and $U_{\calR}$ the uniform distribution on $\calR$.  The source $S$, the predictor $P$
and the distinguisher $\mathsf{D}$ of Figure~\ref{fig:games} are computationally unbounded and
receive the \emph{entire} function table of $H$ as an explicit input.

\section{The conjecture}\label{sec:conjecture}

Notation as in Section~\ref{sec:setting}.  Definition~\ref{def:games} restates the games
formally, for self-containment; the conjecture quantifies over all $\epsilon$-unpredictable
sources.

\begin{figure}[ht]
\centering
\fbox{\begin{minipage}[t]{0.9\linewidth}
\underline{$\Gm\ \Pred^{S}_{P}$}\\[2pt]
$H \sample \Fun(\calK \times \calD, \calR)$;\quad
$(x,z) \sample S(H)$;\quad
$x' \sample P(H,z)$;\quad
$\ret (x = x')$
\end{minipage}}
\par\smallskip
\fbox{\begin{minipage}[t]{0.9\linewidth}
\underline{$\Gm\ \Ext^{S}_{\mathsf{D}}$ (secret seed)}\\[2pt]
$H \sample \Fun(\calK \times \calD, \calR)$;\quad
$(x,z) \sample S(H)$;\quad
$\sd \sample \calK$\\
$y_0 \leftarrow H(\sd,x)$;\quad
$y_1 \sample \calR$;\quad
$b \sample \{0,1\}$;\quad
$b' \sample \mathsf{D}(H, y_b, z)$;\quad
$\ret (b = b')$
\end{minipage}}
\caption{Prediction game (top) and the secret-seed extraction game (bottom); unlike the
public-seed variant (Remark~\ref{rem:pub}), $\mathsf{D}$ does not receive $\sd$.}
\label{fig:games}
\end{figure}

\begin{definition}[the extraction games]\label{def:games}
Set $\Adv^{\pred}_{\calK,\calD,\calR,S}(P) := \Pr[\Pred^{S}_{P} \Rightarrow 1]$ and
\[
  \Adv^{\ext}_{\calK,\calD,\calR}(S,\mathsf{D}) := 2\Pr[\Ext^{S}_{\mathsf{D}} \Rightarrow 1] - 1 .
\]
A source $S$ is \emph{$\epsilon$-unpredictable} if
$\Adv^{\pred}_{\calK,\calD,\calR,S}(P) \le \epsilon$ for every unbounded predictor $P$.
\end{definition}

\begin{conjecture}[secret seed]\label{conj:main}
There is a universal constant $c > 0$, independent of $\calK$, $\calD$, $\calR$ and $\epsilon$,
such that for all nonempty finite $\calK, \calD, \calR$, all $\epsilon \in (0,1]$, every
$\epsilon$-unpredictable source $S$ and every unbounded distinguisher $\mathsf{D}$,
\[
  \Adv^{\ext}_{\calK,\calD,\calR}(S,\mathsf{D})
  \;\le\; \delta(\epsilon, K, D, R)
  \;:=\; c \cdot \sqrt{\frac{\epsilon R + \log_2 D}{K}} \; .
\]
\end{conjecture}

\begin{remark}[the public-seed variant is already resolved]\label{rem:pub}
Giving the seed to $\mathsf{D}$ as well defines a second game, $\Ext\text{-}\mathrm{pub}$, with
advantage $\Adv^{\extpub}$.  Unlike the conjecture above, its natural bound is proved: for the
same universal constant $c = 2$,
\[
  \Adv^{\extpub}_{\calK,\calD,\calR}(S,\mathsf{D})
  \;\le\; c \cdot \Bigl( \sqrt{\epsilon R} + \sqrt{\frac{\log_2 D}{K}} \,\Bigr)
\]
holds unconditionally (see \url{https://crypto-conjura.github.io/c/0004/}).  This is weaker than
Conjecture~\ref{conj:main} whenever the two are compared directly, since
$\sqrt{(\epsilon R + \log_2 D)/K} \le \sqrt{\epsilon R} + \sqrt{\log_2 D/K}$ -- consistent with
hiding the seed only helping the extractor, never hurting it.
\end{remark}

\end{document}
